% ============================================================
% CONCLUSION GÉNÉRALE
% ============================================================
\chapter*{Conclusion Générale}
\addcontentsline{toc}{chapter}{Conclusion Générale}
\markboth{CONCLUSION GÉNÉRALE}{}
\thispagestyle{plain}

\begingroup
\setstretch{1.22}
\setlength{\parskip}{0.55em}

Ce Projet de Fin d'Année avait pour objectif de concevoir et de réaliser
une \textbf{plateforme intelligente de gestion du recrutement} pour le
réseau multi-agences \textbf{DAAM}. Face à la dispersion des informations
et au besoin de coordination entre Administrateur, RH et candidats, la
solution proposée digitalise le parcours complet~: création des offres,
candidatures, présélection, entretiens, embauche, pilotage RH, puis
canal de \textbf{réclamations}.

Sur le plan organisationnel, le projet a suivi une démarche
\textbf{SCRUM}, découpée en sprints successifs. Cette organisation a
permis de livrer progressivement des incréments validables~: socle
d'authentification et d'administration, gestion des offres et zones,
candidatures et QCM, entretiens et tableaux de bord, enfin microservice
de réclamations.

Sur le plan technique, l'architecture retenue combine un frontend
\textbf{Angular} et un backend en \textbf{microservices Spring Boot},
orchestrés par \textbf{Eureka} et exposés via une \textbf{API Gateway}.
La sécurité repose sur \textbf{JWT} et un contrôle d'accès par rôles.
Cette modularité facilite l'évolution du système sans remettre en cause
l'ensemble de la plateforme.

Enfin, l'\textbf{industrialisation DevOps} (Docker, Jenkins, SonarQube,
Docker Hub) a permis d'automatiser la construction, l'analyse qualité et
le déploiement des microservices \textbf{et du frontend Angular}. La
plateforme n'est ainsi plus uniquement un prototype local~: elle dispose
d'une chaîne de livraison reproductible, adaptée à un contexte
professionnel.

Plusieurs \textbf{perspectives} restent ouvertes~: enrichir les
notifications in-app, introduire un fichier \texttt{docker-compose} pour
démarrer l'ensemble de la stack, renforcer la couverture de tests, et
éventuellement conteneuriser MySQL. Ces axes prolongeraient naturellement
le travail réalisé dans le cadre de ce PFA.

En résumé, le projet a permis de relier analyse des besoins, conception
UML, développement microservices et industrialisation DevOps autour d'un
besoin métier concret~: améliorer l'efficacité et la transparence du
recrutement chez DAAM.

\endgroup
\clearpage
